\documentclass[11pt]{article}
\usepackage[margin=1in]{geometry}
\usepackage{natbib}
\usepackage{booktabs}
\usepackage{amsmath}
\usepackage{graphicx}
\usepackage{hyperref}
\usepackage{authblk}

\title{Design Principles for Inflammasome Inhibition by Pyrin-Only-Proteins}

\author[1]{Authors}
\affil[1]{Department of Structural Biology, University of Pittsburgh, Pittsburgh, PA, USA}

\date{}

\begin{document}
\maketitle

\begin{abstract}
Inflammasomes are filamentous signaling platforms essential for host defense against various intracellular calamities such as pathogen invasion and genotoxic stresses. However, dysregulated inflammasomes cause an array of human diseases including autoinflammatory disorders and cancer. It was recently identified that endogenous pyrin-only-proteins (POPs) regulate inflammasomes by directly inhibiting their filament assembly. Here, by combining Rosetta \textit{in silico}, \textit{in vitro}, and \textit{in cellulo} methods, we investigate the target specificity and inhibition mechanisms of POPs. In contrast to a previous report, we find that POP1 is a poor inhibitor of the central inflammasome adaptor ASC. Instead, POP1 inhibits the assembly of upstream receptor PYD filaments such as those of AIM2, IFI16, NLRP3, and NLRP6. Moreover, not only does POP2 directly suppress the nucleation of ASC, but it can also inhibit the elongation of receptor filaments. In addition to inhibiting the elongation of AIM2 and NLRP6 filaments, POP3 potently suppresses the nucleation of ASC. Our Rosetta analyses and biochemical experiments consistently suggest that a combination of favorable and unfavorable interactions between POPs and PYDs is necessary for effective recognition and inhibition. Together, we reveal the intrinsic target redundancy of POPs and their inhibitory mechanisms.
\end{abstract}

\section{Introduction}

Inflammasomes are filamentous signaling platforms integral to host innate defense against a wide range of intracellular catastrophes, including ionizing irradiation, genotoxic chemicals, and pathogen invasion \citep{broz2016,kagan2014,medzhitov2018}. However, persisting inflammasome activities lead to several human maladies including numerous autoinflammatory disorders, cancer, and even severe COVID-19 \citep{karki2017,tartey2020}. Understanding how inflammasome assemblies are regulated at the molecular level can therefore provide key insights into developing strategies for preventing and treating various diseases.

An array of molecular signatures arising from pathogenic conditions induces the oligomerization of inflammasome receptors, resulting in filament assembly by their pyrin domains (PYDs) \citep{lu2014,lu2015b}. The upstream PYD oligomers then induce the filament assembly by the PYD of the central adaptor ASC (ASC$^{\text{PYD}}$), resulting in oligomerization/filamentation of its CARD \citep{lu2014,kagan2014}. It was recently identified that endogenous pyrin-only-proteins (POPs) regulate inflammasomes by directly inhibiting their filament assembly \citep{dealmeida2015,khare2014,ratsimandresy2017}. However, the precise mechanisms by which POPs recognize and inhibit different PYD filaments, and the intrinsic target specificity of individual POPs, have remained poorly understood \citep{devi2020,indramohan2018}.

Here, we combine Rosetta-based \textit{in silico} analyses, \textit{in vitro} biochemical assays (FRET-based polymerization and negative-stain electron microscopy), and \textit{in cellulo} experiments (fluorescence microscopy in HEK293T cells) to systematically dissect the target specificity and inhibition mechanisms of POP1, POP2, and POP3 against five inflammasome PYDs: ASC, AIM2, IFI16, NLRP3, and NLRP6.

\section{Results}

\subsection{Rosetta \textit{in silico} analyses of POP--PYD interactions}

To elucidate how POPs recognize and regulate the assembly of different PYD filaments, we implemented a Rosetta-based \textit{in silico} approach previously developed to define the directionality of the AIM2-ASC inflammasome \citep{matyszewski2021}. PYDs assemble into helical filaments in which each protomer provides six unique protein--protein interaction surfaces (Type~1a/b, 2a/b, and 3a/b) \citep{lu2014,lu2015b}. We created honeycomb-like sideviews of PYD filaments and calculated Rosetta interface energies ($\Delta G$s) for each homotypic PYD filament and for POP$\cdot$PYD interactions by replacing the center protomer with each POP (Figure~1).

The overall $\Delta G$s for individual PYD filaments were significantly more favorable than those from any POP$\cdot$PYD interactions, suggesting that excess POPs would be necessary to inhibit the assembly of inflammasome PYDs (Figure~1A--E). POP1 showed more favorable overall $\Delta G$s for ASC$^{\text{PYD}}$ compared to POP2 or POP3 ($\Delta G = -24.7$ for POP1$\cdot$ASC$^{\text{PYD}}$ vs. $-0.8$ for POP2$\cdot$ASC), seemingly supporting the previous report that POP1 is the master regulator of ASC \citep{dealmeida2015}. For AIM2$^{\text{PYD}}$, all three POPs showed comparable overall $\Delta G$s on the bottom half (Figure~1B). The analyses predicted potential targeting of NLRP3$^{\text{PYD}}$ by POP1 ($\Delta G = -30.5$) and of NLRP6$^{\text{PYD}}$ with similar preference (Figure~1D--E).

\subsection{POP1 does not directly inhibit ASC}

Surprisingly, compared to co-transfecting eGFP alone, POP1-eGFP minimally inhibited the filament assembly of ASC$^{\text{PYD}}$-mCherry ($\leq$20\% suppression; Figure~2A--B). By contrast, POP2-eGFP and POP3-eGFP significantly reduced ASC$^{\text{PYD}}$-mCherry filaments, with POP2 being more effective. POP1 reduced ASC$^{\text{FL}}$ puncta by approximately 30\%, yet POP2 and POP3 were again more effective (up to approximately 80\% reduction; Figure~2C--D).

FRET-based polymerization assays confirmed these observations. POP1 (even at 150~$\mu$M) did not affect ASC$^{\text{PYD}}$ polymerization (Figure~2E). Both POP2 and POP3 significantly prolonged the initial lag phase (up to approximately 15 min delay in nucleation) while moderately affecting elongation ($\leq$20\% reduction in slope; IC$_{50}$s in Figure~2E). nsEM confirmed no ASC$^{\text{PYD}}$ filaments in the presence of POP2, and only a few filaments with POP3 (Figure~2F).

\subsection{A combination of favorable and unfavorable interactions is necessary for inhibition}

Re-examining the Rosetta results in light of biochemical data, we noted that while all three POPs contain favorable interaction surfaces for ASC$^{\text{PYD}}$ ($\Delta\Delta G \leq 3.5$), both POP2 and POP3 show multiple interfaces with significantly unfavorable $\Delta G$s compared to homotypic interactions, whereas POP1$\cdot$ASC$^{\text{PYD}}$ interfaces lack such negative interactions ($\Delta\Delta G \geq 10.0$; Figure~2, Supplement~2). These observations raised the hypothesis that a combination of favorable (recognition) and unfavorable (repulsion) interfaces is necessary for POPs to interfere with PYD filament assembly.

\subsection{Inhibition of ALR assembly by POPs}

Consistent with the amended interpretation, all three POPs showed mixtures of favorable and unfavorable interactions with both AIM2$^{\text{PYD}}$ and IFI16$^{\text{PYD}}$ (Figure~3, Supplement~1). In cells, POP1-eGFP reduced AIM2$^{\text{PYD}}$ and IFI16$^{\text{PYD}}$ filaments more effectively than ASC$^{\text{PYD}}$ ($\sim$60\% vs. $\sim$20\% at 1200~ng POP1; Figure~3B,D). POP2 and POP3 essentially obliterated ALR filament assembly (Figure~3A--D).

FRET assays showed that all three POPs decreased the elongation rate of AIM2$^{\text{PYD}}$ polymerization without significantly affecting nucleation, with POP3 being most effective (Figure~3E). nsEM showed that POP1 produced shorter and fewer AIM2$^{\text{PYD}}$ filaments, while POP2/3 abrogated filament formation entirely (Figure~3F). POP2 and POP3 also inhibited dsDNA binding of AIM2$^{\text{FL}}$, while only POP3 was inhibitory toward IFI16$^{\text{FL}}$ (Figure~3, Supplement~2).

\subsection{POP1 targets upstream NLR receptors}

Our data indicated that POP1 is ineffective against ASC but more effective against AIM2 and IFI16, suggesting POP1 functions as a ``decoy'' ASC targeting upstream receptors. Indeed, POP1 was more effective in suppressing NLRP3$^{\text{PYD}}$ and NLRP6$^{\text{PYD}}$ filament assembly than ASC$^{\text{PYD}}$ (70\% and 60\% suppression, respectively, vs. 20\% for ASC$^{\text{PYD}}$ at 1200~ng POP1; Figure~4A--D).

Using FRET-labeled NLRP6$^{\text{PYD}}$, POP1 predominantly inhibited elongation, POP2 appeared to interfere with nucleation, and POP3 mostly reduced elongation kinetics while also affecting nucleation (Figure~4E). nsEM confirmed that filament number and length were significantly reduced in the presence of POPs (Figure~4F).

\section{Discussion}

Our systematic investigation reveals that the target specificity of POPs is more nuanced and redundant than previously appreciated. In contrast to the prevailing model that POP1 is the primary inhibitor of ASC \citep{dealmeida2015}, we find that POP1 is a poor direct inhibitor of ASC assembly. Instead, POP1 preferentially targets upstream receptor PYDs including AIM2, IFI16, NLRP3, and NLRP6, functioning as a decoy ASC that intercepts upstream signal propagation.

Our combined Rosetta and biochemical analyses reveal a key design principle: effective POP-mediated inhibition requires a combination of both favorable (for recognition/targeting) and unfavorable (for disruption) interaction interfaces between POPs and target PYDs \citep{brandstetter2008}. This mechanism explains why POP1 fails to inhibit ASC despite having favorable overall interaction energies --- it lacks the unfavorable interfaces needed for disruption.

The intrinsic target redundancy of POPs likely serves an important biological purpose. Multiple layers of inhibition --- POP1 targeting upstream receptors, POP2 potently inhibiting both ASC nucleation and receptor elongation, and POP3 combining ASC nucleation suppression with ALR elongation inhibition --- provide robust negative regulation of the inflammasome pathway (Figure~5). This architecture ensures that the inflammasome system, which must balance rapid activation for host defense against pathological hyperactivation, maintains appropriate signaling thresholds \citep{broz2016,kagan2014,indramohan2018}.

\section{Materials and Methods}

\subsection{Rosetta simulation}
The InterfaceAnalyzer script in Rosetta was used to determine the interaction energy at individual interfaces of the honeycomb \citep{matyszewski2021}. We used the cryo-EM structures of ASC$^{\text{PYD}}$ (PDB: 3j63; \citealt{lu2014}), AIM2$^{\text{PYD}}$ (PDB: 7k3r; \citealt{matyszewski2021}), NLRP3$^{\text{PYD}}$ (PDB: 7pdz; \citealt{hochheiser2022a}), and NLRP6$^{\text{PYD}}$ (PDB: 6ncv; \citealt{shen2019}). For IFI16$^{\text{PYD}}$, we used the eGFP-AIM2$^{\text{PYD}}$ filament (PDB: 6mb2; \citealt{lu2015}) as a template. For POPs, we used the crystal structure of POP1 (PDB: 4qob) and generated homology models of POP2 and POP3.

\subsection{Cell culture and imaging}
Each protein was cloned into pCMV6 vector containing C-terminal mCherry (inflammasome PYDs) or eGFP (POPs). HEK293T cells (ATCC, CRL-11268) were seeded into 12-well plates (0.1 $\times$ 10$^6$ per well) with round cover glass. Plasmids were co-transfected using Lipofectamine 2000. After 16 hours, cells were washed, fixed, stained with DAPI, and imaged using fluorescence microscopy ($n \geq 4$ biological replicates).

\subsection{FRET-based polymerization assay}
Donor- and acceptor-labeled PYDs (2.5~$\mu$M total) were mixed with increasing concentrations of POPs, and time-dependent FRET signals were monitored. IC$_{50}$s were calculated from at least three independent measurements. The assay captures two phases: rate-limiting nucleation (initial lag) followed by elongation (exponential/linear phase) \citep{matyszewski2018}.

\subsection{Negative-stain electron microscopy}
PYD filaments (2.5--5~$\mu$M) were incubated with or without POPs for 30 minutes, applied to glow-discharged carbon-coated grids, and stained with 2\% uranyl acetate for imaging \citep{matyszewski2019}.

\subsection{DNA-binding assay}
Fraction of fluorescein-labeled 60-bp dsDNA (10~nM) bound to AIM2$^{\text{FL}}$ (100~nM) or IFI16 (200~nM) was monitored with increasing concentrations of POPs in at least two independent experiments.

\bibliographystyle{plainnat}
\bibliography{references}

\end{document}
